\chapter{2013年江西卷（理）}

\section{单选题}

\begin{problem}
已知集合 \(M = \{1, 2, z\i\}\)，\(\i\) 为虚数单位，\(N = \{3, 4\}\)，\(M \cap N = \{4\}\)，则复数 \(z\) 等于（\quad）

\choices
    {\(-2\i\)}
    {\(2\i\)}
    {\(-4\i\)}
    {\(4\i\)}
\end{problem}

\begin{answer}
C
\end{answer}

\begin{solution}
因为 \(1\)，\(2\notin N\)，而 \(M\cap N=\{4\}\)，所以 \(z\i=4\). 两边同除以 \(\i\)，得
\[
z=\frac{4}{\i}=-4\i
\]
故选 C.
\end{solution}

\begin{problem}
函数 \( y = \sqrt{x} \ln (1 - x) \) 的定义域为（\quad）

\choices
    {\((0,1)\)}
    {\([0,1)\)}
    {\((0,1]\)}
    {\([0,1]\)}
\end{problem}

\begin{answer}
B
\end{answer}

\begin{solution}
根式有意义要求 \(x\geqslant0\)，对数有意义要求 \(1-x>0\)，即 \(x<1\). 两条件同时成立时
\[
x\in[0,1)
\]
故选 B.
\end{solution}

\begin{problem}
等比数列 \(x\)，\(3x + 3\)，\(6x + 6\)，\(\cdots\) 的第四项等于（\quad）

\choices
    {\(-24\)}
    {0}
    {12}
    {24}
\end{problem}

\begin{answer}
A
\end{answer}

\begin{solution}
等比数列相邻两项的比必须有意义，所以 \(x\ne0\)，且 \(3x+3\ne0\). 由前两次比值相等，
\[
\frac{3x+3}{x}=\frac{6x+6}{3x+3}
\]
交叉相乘并整理：
\[
(3x+3)^2=x(6x+6)
\Longleftrightarrow
3(x+1)(x+3)=0
\]
候选值为 \(x=-1\)，\(-3\). 当 \(x=-1\) 时第二项、第三项均为 \(0\)，相邻项的比无意义，不能构成题设等比数列；故 \(x=-3\). 此时数列为
\(-3\)，\(-6\)，\(-12\)，\(-24\).
故第四项为 \(-24\)，选 A.
\end{solution}

\begin{problem}
总体由编号为 01，02，…，19，20 的 20 个个体组成. 利用下面的随机数表选取 5 个个体，选取方法从随机数表第 1 行的第 5 列和第 6 列数字开始，由左到右依次选取两个数字，则选出来的第 5 个个体的编号为（\quad）

7816 6572 0802 6314 0702 4369 9728 0198 3204 9234 4935 8200 3623 4869 6938 7481

\choices
    {08}
    {07}
    {02}
    {01}
\end{problem}

\begin{answer}
D
\end{answer}

\begin{solution}
从第 5、6 列开始，每次取相邻两位，得到
\(65\)，\(\ 72\)，\(\ 08\)，\(\ 02\)，\(\ 63\)，\(\ 14\)，\(\ 07\)，\(\ 02\)，\(\ 43\)，\(\ 69\)，\(\ 97\)，\(\ 28\)，\(\ 01\)，\(\ldots\)
编号必须在 \(01\) 到 \(20\) 之间，且已经选过的编号不能重复. 因此依次保留
\(08\)，\(02\)，\(14\)，\(07\)，\(01\).
第五个个体的编号为 \(01\)，故选 D.
\end{solution}

\begin{problem}
\(\left(x^{2} - \frac{2}{x^{3}}\right)^{5}\) 展开式中的常数项为（\quad）

\choices
    {80}
    {\(-80\)}
    {40}
    {\(-40\)}
\end{problem}

\begin{answer}
C
\end{answer}

\begin{solution}
从展开式中取 \(k\) 个 \(x^2\)，取 \(5-k\) 个 \(-2x^{-3}\)，所得项为
\[
\mathrm C_5^k(x^2)^k\left(-2x^{-3}\right)^{5-k}
=\mathrm C_5^k(-2)^{5-k}x^{2k-3(5-k)}
=\mathrm C_5^k(-2)^{5-k}x^{5k-15}
\]
成为常数项必须满足 \(5k-15=0\)，所以 \(k=3\). 常数项系数为
\[
\mathrm C_5^3(-2)^2=10\cdot4=40
\]
故选 C.
\end{solution}

\begin{problem}
若 \(S_{1} = \int_{1}^{2} x^{2} \mathrm{d}x\)，\(S_{2} = \int_{1}^{2} \frac{1}{x} \mathrm{d}x\)，\(S_{3} = \int_{1}^{2} \e^{x} \mathrm{d}x\)，则 \(S_{1}\)，\(S_{2}\)，\(S_{3}\) 的大小关系为（\quad）

\choices
    {\(S_{1} < S_{2} < S_{3}\)}
    {\(S_{2} < S_{1} < S_{3}\)}
    {\(S_{2} < S_{3} < S_{1}\)}
    {\(S_{3} < S_{2} < S_{1}\)}
\end{problem}

\begin{answer}
B
\end{answer}

\begin{solution}
分别计算三个积分：
\(S_1=\left.\frac{x^3}{3}\right|_1^2=\frac{7}{3}\)，\(S_2=\left.\ln x\right|_1^2=\ln2\)
\[
S_3=\left.\e^x\right|_1^2=\e^2-\e=\e(\e-1)
\]
在 \(1<x<2\) 时有 \(0<1/x<1\)，故 \(0<S_2<1<S_1\). 又由指数函数的展开式
\(\e^x=1+x+\frac{x^2}{2}+\cdots\)，在 \(1\leqslant x\leqslant2\) 时有
\(\e^x\geqslant1+x+\frac{x^2}{2}\)，\(1+x+\frac{x^2}{2}-x^2=1+x\left(1-\frac{x}{2}\right)\geqslant1>0\)
故 \(\e^x>x^2\)，从而 \(S_3>S_1\). 因此
\[
S_2<S_1<S_3
\]
故选 B.
\end{solution}

\begin{problem}
阅读如图所示程序框图，如果输出 \(i = 5\)，那么在空白矩形框中应填入的语句为（\quad）

\bitmapfigure[max width=5.8cm]{img/2013/jiangxi_science/q7_fig1.png}

\choices
    {\(S = 2 * i - 2\)}
    {\(S = 2 * i - 1\)}
    {\(S = 2 * i\)}
    {\(S = 2 * i + 4\)}
\end{problem}

\begin{answer}
C
\end{answer}

\begin{solution}
程序开始时 \(i=1\)，\(S=0\)，每次循环先令 \(i\leftarrow i+1\). 当 \(i\) 为偶数时，框图中的语句令 \(S=2i+1\)；当 \(i\) 为奇数时执行空白语句.

若空白语句为 \(S=2i\)，逐次计算得到：当 \(i=2\)，\(3\)，\(4\)，\(5\) 时，
\(S\) 分别为 \(5\)，\(6\)，\(9\)，\(10\).
在 \(i=2\)，\(3\)，\(4\) 时均有 \(S<10\)，到 \(i=5\) 时 \(S=10\)，程序输出 \(i=5\)，符合题意. 其余三个选项分别使程序在 \(i=6\)，\(6\)，\(3\) 时输出，均不符合题意. 因此空白语句为 \(S=2*i\)，选 C.
\end{solution}

\begin{problem}
如图，正方体的底面与正四面体的底面在同一平面 \(\alpha\) 上，且 \(AB \parallel CD\)，正方体的六个面所在的平面与直线 \(CE\)，\(EF\) 相交的平面个数分别记为 \(m\)，\(n\)，那么 \(m + n =\)（\quad）

\bitmapfigure[max width=0.48\linewidth]{img/2013/jiangxi_science/q8_fig1.png}

\choices
    {8}
    {9}
    {10}
    {11}
\end{problem}

\begin{answer}
A
\end{answer}

\begin{solution}
正方体的两个水平面分别为底面 \(\alpha\) 和顶面，另有四个竖直侧面.

直线 \(CE\) 在 \(\alpha\) 内，且正三角形 \(CDE\) 中 \(\angle DCE=60^{\circ}\). 由于 \(CD\parallel AB\)，\(CE\) 不平行于正方体的任一组侧棱方向，所以它与四个竖直侧面所在平面均相交；它包含在底面 \(\alpha\) 内，并与顶面平行，故 \(m=4\).

设 \(F\) 在 \(\alpha\) 上的射影为 \(H\). 正四面体的高线 \(FH\) 通过底面中心，因此 \(EH\perp CD\parallel AB\). 直线 \(EF\) 的水平投影垂直于 \(AB\)，所以它与其中两个竖直侧面所在平面平行，与另外两个竖直侧面所在平面相交；同时它与两个水平面相交，故 \(n=4\).

所以
\[
m+n=4+4=8
\]
故选 A.
\end{solution}

\begin{problem}
过点 \((\sqrt{2}, 0)\) 引直线 \(l\) 与曲线 \(y = \sqrt{1 - x^{2}}\) 相交于 \(A\)，\(B\) 两点，\(O\) 为坐标原点，当 \(\triangle AOB\) 的面积取最大值时，直线 \(l\) 的斜率等于（\quad）

\choices
    {\(\frac{\sqrt{3}}{3}\)}
    {\(-\frac{\sqrt{3}}{3}\)}
    {\(\pm \frac{\sqrt{3}}{3}\)}
    {\(-\sqrt{3}\)}
\end{problem}

\begin{answer}
B
\end{answer}

\begin{solution}
曲线是单位圆的上半圆. 设直线 \(l\) 与原点的距离为 \(d\)，则弦 \(AB\) 的长为
\[
AB=2\sqrt{1-d^2}
\]
以 \(AB\) 为底，原点到 \(AB\) 的距离为高，所以
\[
[\triangle AOB]=\frac12 AB\cdot d=d\sqrt{1-d^2}
\]
令 \(u=d^2\)，则
\[
[\triangle AOB]^2=u(1-u)\leqslant\frac14
\]
等号当且仅当 \(d=1/\sqrt2\).

设直线斜率为 \(k\)，其方程为 \(y=k(x-\sqrt2)\)，于是
\[
d=\frac{|k|\sqrt2}{\sqrt{k^2+1}}
\]
代入 \(d^2=1/2\)：
\[
\frac{2k^2}{k^2+1}=\frac12
\Longrightarrow
3k^2=1
\Longrightarrow
k=\pm\frac{\sqrt3}{3}
\]
当 \(k>0\) 时，单位上半圆的横坐标满足 \(x\leqslant1<\sqrt2\)，从而 \(y=k(x-\sqrt2)<0\)，不能与上半圆有两个交点；故只能取
\[
k=-\frac{\sqrt3}{3}
\]
故选 B.
\end{solution}

\begin{problem}
如图，半径为 1 的半圆 \(O\) 与等边三角形 \(ABC\) 夹在两平行线 \(l_{1}\)，\(l_{2}\) 之间，\(l \parallel l_{1}\)，\(l\) 与半圆相交于 \(F\)，\(G\) 两点，与三角形 \(ABC\) 两边相交于 \(E\)，\(D\) 两点. 设弧 \(FG\) 的长为 \(x\)（\(0 < x < \pi\)），\(y = EB + BC + CD\)，若 \(l\) 从 \(l_{1}\) 平行移动到 \(l_{2}\)，则函数 \(y = f(x)\) 的图象大致是（\quad）

\FigureLayoutDeclare{img/2013/jiangxi_science/q9_fig1.png}{8.0cm}{60.992bp 0bp 44.994bp 0bp}
\bitmapfigure[max width=0.60\linewidth]{img/2013/jiangxi_science/q9_fig1.png}

\choices
    {\choicebitmap[width=0.15\paperwidth]{img/2013/jiangxi_science/q10_opt_a.png}}
    {\choicebitmap[width=0.15\paperwidth]{img/2013/jiangxi_science/q10_opt_b.png}}
    {\choicebitmap[width=0.15\paperwidth]{img/2013/jiangxi_science/q10_opt_c.png}}
    {\choicebitmap[width=0.15\paperwidth]{img/2013/jiangxi_science/q10_opt_d.png}}
\end{problem}

\begin{answer}
D
\end{answer}

\begin{solution}
取 \(l_1\) 为半圆最低处所在的水平线，令 \(t\) 为直线 \(l\) 到 \(l_1\) 的距离. 两条平行线间距为半径 \(1\)，所以等边三角形的高也为 \(1\)，其边长为
\[
s=\frac{2}{\sqrt3}=\frac{2\sqrt3}{3}
\]
由相似三角形，直线 \(l\) 上移的高度为 \(t\) 时，
\(EB=CD=\frac{2t}{\sqrt3}\)，\(y=s+\frac{4t}{\sqrt3}\).
半径为 \(1\) 的下半圆中，若弧 \(FG\) 的长为 \(x\)，则
\[
t=1-\cos\frac{x}{2}
\]
因此
\(f(x)=\frac{2\sqrt3}{3} +\frac{4}{\sqrt3}\left(1-\cos\frac{x}{2}\right)\)，\(0<x<\pi\).
于是
\(f'(x)=\frac{2}{\sqrt3}\sin\frac{x}{2}>0\)，\(f''(x)=\frac{1}{\sqrt3}\cos\frac{x}{2}>0\).
函数递增且图象向上凸，故选 D.
\end{solution}

\section{填空题}

\begin{problem}
函数 \(y = \sin 2x + 2\sqrt{3}\sin^{2}x\) 的最小正周期 \(T\) 为\fillinblank{}.
\end{problem}

\begin{answer}
\(\pi\)
\end{answer}

\begin{solution}
利用 \(\sin^2x=(1-\cos2x)/2\)，得
\[
y=\sin2x+\sqrt3(1-\cos2x)
=\sqrt3+\sin2x-\sqrt3\cos2x
\]
后两项是角频率为 \(2\) 的非零正弦型函数，故最小正周期为
\[
T=\frac{2\pi}{2}=\pi
\]
\end{solution}

\begin{problem}
\(\bs{e}_{1}\)，\(\bs{e}_{2}\) 为单位向量，且 \(\bs{e}_{1}\)，\(\bs{e}_{2}\) 的夹角为 \(\frac{\pi}{3}\)，若 \(\bs{a} = \bs{e}_{1} + 3\bs{e}_{2}\)，\(\bs{b} = 2\bs{e}_{1}\)，则向量 \(\bs{a}\) 在 \(\bs{b}\) 方向上的射影为\fillinblank{}.
\end{problem}

\begin{answer}
\(\frac{5}{2}\)
\end{answer}

\begin{solution}
向量 \(\bs a\) 在 \(\bs b\) 方向上的射影为
\[
\frac{\bs a\cdot\bs b}{|\bs b|}
\]
由 \(|\bs e_1|=1\)，且 \(\angle(\bs e_1,\bs e_2)=\pi/3\)，有
\[
\bs e_1\cdot\bs e_2=\cos\frac{\pi}{3}=\frac12
\]
因此
\(\bs a\cdot\bs b =(\bs e_1+3\bs e_2)\cdot2\bs e_1 =2+6\cdot\frac12=5\)，\(|\bs b|=2\).
所求射影为 \(5/2\).
\end{solution}

\begin{problem}
设函数 \(f(x)\) 在 \((0, +\infty)\) 内可导，且 \(f(\e^x) = x + \e^x\)，则 \(f'(1) =\)\fillinblank{}.
\end{problem}

\begin{answer}
\(2\)
\end{answer}

\begin{solution}
对恒等式 \(f(\e^x)=x+\e^x\) 两边关于 \(x\) 求导，得
\[
\e^x f'(\e^x)=1+\e^x
\]
令 \(x=0\)，则 \(\e^x=1\)，从而
\[
f'(1)=1+1=2
\]
\end{solution}

\begin{problem}
抛物线 \(x^{2} = 2py\)（\(p > 0\)）的焦点为 \(F\)，其准线与双曲线 \(\frac{x^{2}}{3} - \frac{y^{2}}{3} = 1\) 相交于 \(A\)，\(B\) 两点，若 \(\triangle ABF\) 为等边三角形，则 \(p =\)\fillinblank{}.
\end{problem}

\begin{answer}
\(6\)
\end{answer}

\begin{solution}
抛物线 \(x^2=2py\) 的焦点为 \(F(0,p/2)\)，准线为 \(y=-p/2\). 令 \(y=-p/2\) 代入双曲线，得
\[
\frac{x^2}{3}-\frac{p^2}{12}=1
\Longrightarrow
x^2=3+\frac{p^2}{4}
\]
所以可设
\(A\left(-\sqrt{3+\frac{p^2}{4}},-\frac p2\right)\)，\(B\left(\sqrt{3+\frac{p^2}{4}},-\frac p2\right)\).
于是
\[
AB^2=4\left(3+\frac{p^2}{4}\right)=12+p^2
\]
而
\[
AF^2=3+\frac{p^2}{4}+p^2=3+\frac{5p^2}{4}
\]
等边三角形要求 \(AF=AB\)，故
\[
3+\frac{5p^2}{4}=12+p^2
\Longrightarrow
\frac{p^2}{4}=9
\]
由于 \(p>0\)，得 \(p=6\).
\end{solution}

\begin{problem}
设曲线 \(C\) 的参数方程为 \(\begin{cases}
x = t,\\
y = t^2
\end{cases}\)（\(t\) 为参数），若以直角坐标系的原点为极点，\(x\) 轴的正半轴为极轴建立极坐标系，则曲线 \(C\) 的极坐标方程为\fillinblank{}.
\end{problem}

\begin{answer}
\(\rho\cos^2\theta - \sin\theta = 0\)
\end{answer}

\begin{solution}
极坐标与直角坐标的关系为
\(x=\rho\cos\theta\)，\(y=\rho\sin\theta\).
参数方程消去参数得到 \(y=x^2\)，代入极坐标关系：
\[
\rho\sin\theta=(\rho\cos\theta)^2
\]
对非原点点可化为
\[
\rho\cos^2\theta-\sin\theta=0
\]
原点可取 \(\rho=0\)，\(\theta=0\) 表示，故所求方程为
\[
\rho\cos^2\theta-\sin\theta=0
\]
\end{solution}

\begin{problem}
在实数范围内，不等式 \(||x - 2| - 1| \leqslant 1\) 的解集为\fillinblank{}.
\end{problem}

\begin{answer}
\([0,4]\)
\end{answer}

\begin{solution}
令 \(u=|x-2|\)，则 \(u\geqslant0\). 原不等式化为
\[
|u-1|\leqslant1
\Longleftrightarrow
-1\leqslant u-1\leqslant1
\Longleftrightarrow
0\leqslant u\leqslant2
\]
因此 \(|x-2|\leqslant2\)，即
\[
0\leqslant x\leqslant4
\]
故解集为 \([0,4]\).
\end{solution}

\section{解答题}

\begin{problem}
在 \(\triangle ABC\) 中，角 \(A\)，\(B\)，\(C\) 所对的边分别为 \(a\)，\(b\)，\(c\). 已知 \(\cos C + (\cos A - \sqrt{3}\sin A)\cos B = 0\).

\begin{enumerate}
\item 求角 \(B\) 的大小；
\item 若 \(a + c = 1\)，求 \(b\) 的取值范围.
\end{enumerate}
\end{problem}

\begin{solution}
\begin{enumerate}
\item 因为 \(A+B+C=\pi\)，所以
\[
\cos C=-\cos(A+B)=-\cos A\cos B+\sin A\sin B
\]
代入已知条件：
\[
0=-\cos A\cos B+\sin A\sin B
+(\cos A-\sqrt3\sin A)\cos B
\]
整理得
\[
0=\sin A(\sin B-\sqrt3\cos B)
\]
三角形内角 \(A\in(0,\pi)\)，故 \(\sin A>0\)，从而
\[
\sin B-\sqrt3\cos B=0
\]
又 \(B\in(0,\pi)\)，只能有 \(\tan B=\sqrt3\)，所以
\[
B=\frac{\pi}{3}
\]

\item 由余弦定理及 \(\cos B=1/2\)，
\[
b^2=a^2+c^2-2ac\cos B
=a^2+c^2-ac
=(a+c)^2-3ac
=1-3ac
\]
因为 \(a\)，\(c>0\) 且 \(a+c=1\)，有
\[
0<ac\leqslant\frac14
\]
因此
\[
\frac14\leqslant b^2<1
\]
从而
\[
\frac12\leqslant b<1
\]
端点 \(b=1/2\) 在 \(a=c=1/2\) 时可以取得；\(b=1\) 需要 \(ac=0\)，不符合三角形边长为正，故
\[
b\in\left[\frac12,1\right)
\]
\end{enumerate}
\end{solution}

\begin{problem}
正项数列 \(\{a_{n}\}\) 的前 \(n\) 项和 \(S_{n}\) 满足 \(S_{n}^{2} - (n^{2} + n - 1)S_{n} - (n^{2} + n) = 0\).

\begin{enumerate}
\item 求数列 \(\{a_{n}\}\) 的通项公式；
\item 令 \(b_{n} = \frac{n + 1}{(n + 2)^2 a_{n}^{2}}\)，数列 \(\{b_{n}\}\) 的前 \(n\) 项和为 \(T_{n}\). 证明：对于任意 \(n \in \N^{*}\)，都有 \(T_{n} < \frac{5}{64}\).
\end{enumerate}
\end{problem}

\begin{solution}
\begin{enumerate}
\item 记 \(m=n^2+n\)，则题设方程为
\[
S_n^2-(m-1)S_n-m=(S_n-m)(S_n+1)=0
\]
由于数列各项为正，\(S_n>0\)，故
\[
S_n=n^2+n
\]
当 \(n=1\) 时 \(a_1=S_1=2\)；当 \(n\geqslant2\) 时，
\[
a_n=S_n-S_{n-1}
=(n^2+n)-\bigl((n-1)^2+(n-1)\bigr)=2n
\]
所以
\(a_n=2n\)，\((n\in\N^*)\).

\item 代入 \(a_n=2n\)：
\[
b_n=\frac{n+1}{4n^2(n+2)^2}
\]
又
\[
\frac1{n^2}-\frac1{(n+2)^2}
=\frac{4(n+1)}{n^2(n+2)^2}
\]
故
\[
b_n=\frac1{16}\left(\frac1{n^2}-\frac1{(n+2)^2}\right)
\]
于是
\[
T_n=\frac1{16}\sum_{k=1}^n
\left(\frac1{k^2}-\frac1{(k+2)^2}\right)
=\frac1{16}\left(1+\frac14-\frac1{(n+1)^2}-\frac1{(n+2)^2}\right)
\]
因此
\[
T_n=\frac5{64}-\frac1{16}\left(\frac1{(n+1)^2}+\frac1{(n+2)^2}\right)
<\frac5{64}
\]
证毕.
\end{enumerate}
\end{solution}

\begin{problem}
小波以游戏方式决定是参加学校合唱团还是参加学校排球队，游戏规则为：以 \(O\) 为起点，再从 \(A_{1}\)，\(A_{2}\)，\(A_{3}\)，\(A_{4}\)，\(A_{5}\)，\(A_{6}\)，\(A_{7}\)，\(A_{8}\)（如图）这 8 个点中任取两点，分别为终点得到两个向量，记这两个向量的数量积为 \(X\). 若 \(X = 0\) 就参加学校合唱团，否则就参加学校排球队.

\begin{enumerate}
\item 求小波参加学校合唱团的概率；
\item 求 \(X\) 的分布列和数学期望.
\end{enumerate}

\bitmapfigure[max width=0.36\linewidth]{img/2013/jiangxi_science/q19_fig1.png}
\end{problem}

\begin{solution}
由图可知 \(\overrightarrow{OA_1}=(1,0)\)，\(\overrightarrow{OA_2}=(1,1)\)，\(\overrightarrow{OA_3}=(0,1)\)，\(\overrightarrow{OA_4}=(-1,1)\)；\(\overrightarrow{OA_5}=(-1,0)\)，\(\overrightarrow{OA_6}=(-1,-1)\)，\(\overrightarrow{OA_7}=(0,-1)\)，\(\overrightarrow{OA_8}=(1,-1)\).

从 8 个点中任取两个点，共有 \(\mathrm C_8^2=28\) 个等可能的点对. 逐一计算两向量的数量积，可按点对计数：
当 \(X=-2\) 时，点对为 \((A_2,A_6)\)、\((A_4,A_8)\)，共 \(2\) 对.
当 \(X=-1\) 时，点对为 \((A_1,A_4)\)、\((A_1,A_5)\)、\((A_1,A_6)\)、\((A_2,A_5)\)、\((A_2,A_7)\)、\((A_3,A_6)\)、\((A_3,A_7)\)、\((A_3,A_8)\)、\((A_4,A_7)\)、\((A_5,A_8)\)，共 \(10\) 对.
当 \(X=0\) 时，点对为 \((A_1,A_3)\)、\((A_1,A_7)\)、\((A_2,A_4)\)、\((A_2,A_8)\)、\((A_3,A_5)\)、\((A_4,A_6)\)、\((A_5,A_7)\)、\((A_6,A_8)\)，共 \(8\) 对.
当 \(X=1\) 时，点对为 \((A_1,A_2)\)、\((A_1,A_8)\)、\((A_2,A_3)\)、\((A_3,A_4)\)、\((A_4,A_5)\)、\((A_5,A_6)\)、\((A_6,A_7)\)、\((A_7,A_8)\)，共 \(8\) 对.

\begin{enumerate}
\item
\[
P(X=0)=\frac8{28}=\frac27
\]

\item
\(P(X=-2)=\frac1{14}\)，\(P(X=-1)=\frac5{14}\)，\(P(X=0)=\frac27\)，\(P(X=1)=\frac27\).
因此
\[
E(X)=(-2)\frac1{14}+(-1)\frac5{14}
+0\cdot\frac27+1\cdot\frac27
=-\frac3{14}
\]
\end{enumerate}
\end{solution}

\begin{problem}
如图，四棱锥 \(P - ABCD\) 中，\(PA \perp\) 平面 \(ABCD\)，\(E\) 为 \(BD\) 的中点，\(G\) 为 \(PD\) 的中点，\(\triangle DAB \cong \triangle DCB\)，\(EA = EB = AB = 1\)，\(PA = \frac{3}{2}\)，连接 \(CE\) 并延长交 \(AD\) 于 \(F\).

\begin{enumerate}
\item 求证：\(AD \perp\) 平面 \(CFG\)；
\item 求平面 \(BCP\) 与平面 \(DCP\) 的夹角的余弦值.
\end{enumerate}

\bitmapfigure[max width=0.46\linewidth]{img/2013/jiangxi_science/q20_fig1.png}
\end{problem}

\begin{solution}
在底面内建立平面直角坐标系，取
\(A=(0,0,0)\)，\(B=(1,0,0)\).
因为 \(E\) 是 \(BD\) 的中点，且 \(EA=EB=AB=1\)，三角形 \(AEB\) 为等边三角形. 取
\(E=\left(\frac12,\frac{\sqrt3}{2},0\right)\)，\(D=2E-B=(0,\sqrt3,0)\).
由 \(\triangle DAB\cong\triangle DCB\)，对应边给出 \(DC=DA=\sqrt3\)，\(CB=BA=1\). 取图示位置可得
\[
C=\left(\frac32,\frac{\sqrt3}{2},0\right)
\]
直线 \(CE\) 为 \(y=\sqrt3/2\)，直线 \(AD\) 为 \(x=0\)，故
\[
F=\left(0,\frac{\sqrt3}{2},0\right)
\]
又 \(PA\perp\) 平面 \(ABCD\)，\(PA=3/2\)，所以可取
\(P=\left(0,0,\frac32\right)\)，\(G=\left(0,\frac{\sqrt3}{2},\frac34\right)\).

\begin{enumerate}
\item
\(\overrightarrow{CF}=\left(-\frac32,0,0\right)\)，\(\overrightarrow{FG}=\left(0,0,\frac34\right)\)，\(\overrightarrow{AD}=(0,\sqrt3,0)\).
因此
\(\overrightarrow{AD}\cdot\overrightarrow{CF}=0\)，\(\overrightarrow{AD}\cdot\overrightarrow{FG}=0\).
直线 \(CF\) 与 \(FG\) 相交于 \(F\)，且都在平面 \(CFG\) 内，所以
\[
AD\perp\text{平面 }CFG
\]

\item 取平面 \(BCP\) 的两个方向向量
\(\overrightarrow{BC}=\left(\frac12,\frac{\sqrt3}{2},0\right)\)，\(\overrightarrow{BP}=\left(-1,0,\frac32\right)\)
可取其法向量
\[
\bs n_1=\overrightarrow{BC}\times\overrightarrow{BP}
=\left(\frac{3\sqrt3}{4},-\frac34,\frac{\sqrt3}{2}\right)
\]
同理，平面 \(DCP\) 的两个方向向量可取
\(\overrightarrow{DC}=\left(\frac32,-\frac{\sqrt3}{2},0\right)\)，\(\overrightarrow{DP}=\left(0,-\sqrt3,\frac32\right)\)
可取其法向量
\[
\bs n_2=\overrightarrow{DC}\times\overrightarrow{DP}
=\left(-\frac{3\sqrt3}{4},-\frac94,-\frac{3\sqrt3}{2}\right)
\]
取其相反方向的法向量
\[
\bs n_2=\left(\frac{3\sqrt3}{4},\frac94,\frac{3\sqrt3}{2}\right)
\]
故两平面的夹角余弦为
\[
\frac{|\bs n_1\cdot\bs n_2|}{|\bs n_1|\,|\bs n_2|}
=\frac{9/4}{\sqrt3\cdot\frac{3\sqrt6}{2}}
=\frac{\sqrt2}{4}
\]
\end{enumerate}
\end{solution}

\begin{problem}
如图，椭圆 \(C: \frac{x^2}{a^2} + \frac{y^2}{b^2} = 1\)（\(a > b > 0\)）经过点 \(P\left(1, \frac{3}{2}\right)\)，离心率 \(e = \frac{1}{2}\)，直线 \(l\) 的方程为 \(x = 4\).

\begin{enumerate}
\item 求椭圆 \(C\) 的方程；
\item \(AB\) 是经过右焦点 \(F\) 的任一弦（不经过点 \(P\)），设直线 \(AB\) 与直线 \(l\) 相交于点 \(M\)，记 \(PA\)，\(PB\)，\(PM\) 的斜率分别为 \(k_{1}\)，\(k_{2}\)，\(k_{3}\). 问：是否存在常数 \(\lambda\)，使得 \(k_{1} + k_{2} = \lambda k_{3}\)？若存在，求 \(\lambda\) 的值；若不存在，请说明理由.
\end{enumerate}

\bitmapfigure[max width=0.42\linewidth]{img/2013/jiangxi_science/q21_fig1.png}
\end{problem}

\begin{solution}
\begin{enumerate}
\item 由离心率条件
\(\frac ca=\frac12\)，\(c^2=a^2-b^2\)
得 \(c=a/2\)，从而 \(b^2=3a^2/4\). 将点 \(P(1,3/2)\) 代入椭圆方程：
\[
\frac1{a^2}+\frac{(3/2)^2}{3a^2/4}=1
\Longrightarrow
\frac1{a^2}+\frac3{a^2}=1
\]
所以 \(a^2=4\)，\(b^2=3\)，椭圆方程为
\[
\frac{x^2}{4}+\frac{y^2}{3}=1
\]

\item 右焦点为 \(F=(1,0)\). 直线 \(AB\) 不可能为竖直线，因为竖直线 \(x=1\) 不与 \(x=4\) 相交于有限点 \(M\). 设
\[
AB:\ y=k(x-1)
\]
令 \(t=x-1\)，将 \(x=t+1\) 和 \(y=kt\) 代入椭圆方程，得
\[
\frac{(t+1)^2}{4}+\frac{k^2t^2}{3}=1
\Longleftrightarrow
(3+4k^2)t^2+6t-9=0
\]
设交点 \(A\)，\(B\) 对应的根为 \(t_1\)，\(t_2\)，则
\(t_1+t_2=-\frac6{3+4k^2}\)，\(t_1t_2=-\frac9{3+4k^2}\)
从而
\[
\frac1{t_1}+\frac1{t_2}
=\frac{t_1+t_2}{t_1t_2}
=\frac23
\]
由于 \(P=(1,3/2)\)，
\(k_1=k-\frac{3}{2t_1}\)，\(k_2=k-\frac{3}{2t_2}\)
所以
\[
k_1+k_2
=2k-\frac32\left(\frac1{t_1}+\frac1{t_2}\right)
=2k-1
\]
直线 \(AB\) 与 \(x=4\) 交于
\[
M=(4,3k)
\]
故
\[
k_3=\frac{3k-3/2}{4-1}=k-\frac12
\]
因此
\[
k_1+k_2=2k-1=2\left(k-\frac12\right)=2k_3
\]
存在所求常数，且
\[
\lambda=2
\]
\end{enumerate}
\end{solution}

\begin{problem}
已知函数 \(f(x) = a\left(1 - 2\left|x - \frac{1}{2}\right|\right)\)，\(a\) 为常数且 \(a > 0\).

\begin{enumerate}
\item 证明：函数 \(f(x)\) 的图象关于直线 \(x = \frac{1}{2}\) 对称；
\item 若 \(x_0\) 满足 \(f(f(x_0)) = x_0\)，且 \(f(x_0) \neq x_0\)，则称 \(x_0\) 为函数 \(f(x)\) 的二阶周期点. 如果 \(f(x)\) 有两个二阶周期点 \(x_1\)，\(x_2\)，试确定 \(a\) 的取值范围；
\item 对于第 2 问中的 \(x_1\)，\(x_2\) 和 \(a\)，设 \(x_3\) 为函数 \(f(f(x))\) 的最大值点，\(A(x_1, f(f(x_1)))\)，\(B(x_2, f(f(x_2)))\)，\(C(x_3, 0)\)，记 \(\triangle ABC\) 的面积为 \(S(a)\)，讨论 \(S(a)\) 的单调性.
\end{enumerate}
\end{problem}

\begin{solution}
\begin{enumerate}
\item 对任意 \(x\)，有
\[
f(1-x)
=a\left(1-2\left|1-x-\frac12\right|\right)
=a\left(1-2\left|x-\frac12\right|\right)
=f(x)
\]
因此函数图象关于直线 \(x=1/2\) 对称.

\item 将函数写成分段形式：
\[
f(x)=
\begin{cases}
2ax,&x\leqslant\frac12,\\
2a(1-x),&x\geqslant\frac12.
\end{cases}
\]
设一个非平凡二循环的两个点为 \(u<v\)，即 \(f(u)=v\)，\(f(v)=u\). 若 \(u\)，\(v\leqslant1/2\)，由于 \(f\) 在 \((-\infty,1/2]\) 上递增，有 \(v=f(u)\leqslant f(v)=u\)，与 \(u<v\) 矛盾. 若 \(u\)，\(v\geqslant1/2\)，则
\(v=2a(1-u)\)，\(u=2a(1-v)\).
两式相减得 \((1-2a)(v-u)=0\). 因 \(u\ne v\)，只能有 \(a=1/2\)，此时两式又给出 \(u+v=1\)，不可能有两个不同的点都大于等于 \(1/2\). 所以非平凡二循环必满足 \(u<1/2<v\).

于是
\(x_2=f(x_1)=2ax_1\)，\(x_1=f(x_2)=2a(1-x_2)\).
解得
\(x_1=\frac{2a}{1+4a^2}\)，\(x_2=\frac{4a^2}{1+4a^2}\).
当 \(a\ne1/2\) 时有 \(x_1<1/2\)；并且 \(x_2>1/2\) 当且仅当 \(a>1/2\). 在 \(a=1/2\) 时两点重合，不能构成非平凡二循环；在 \(0<a<1/2\) 时 \(x_2<1/2\)，不满足二循环两点分居 \(1/2\) 两侧的条件，也不存在这样的二循环. 因此
\[
a\in\left(\frac12,+\infty\right)
\]

\item 对 \(a>1/2\)，函数 \(f\) 的最大值为 \(a\)，所以 \(f(f(x))\leqslant a\). 等号当且仅当 \(f(x)=1/2\)，其两个最大值点为
\(\frac1{4a}\)，\(1-\frac1{4a}\).
由于 \(f(f(x_1))=x_1\)，\(f(f(x_2))=x_2\)，故
\(A=(x_1,x_1)\)，\(B=(x_2,x_2)\).
题面把最大值点记为单个 \(x_3\)，但这里实际有两个最大值点，故需分别讨论.

若取较小的最大值点 \(x_3=1/(4a)\)，则三点构成的面积为
\(C=\left(\frac1{4a},0\right)\)，\(S(a)=\frac12(x_2-x_1)x_3 =\frac12\cdot\frac{2a(2a-1)}{1+4a^2}\cdot\frac1{4a} =\frac{2a-1}{4(1+4a^2)}\).
求导：
\[
S'(a)
=\frac{2+8a-8a^2}{4(1+4a^2)^2}
=\frac{-2(4a^2-4a-1)}{4(1+4a^2)^2}
\]
在 \(a>1/2\) 时，临界点为
\[
a_0=\frac{1+\sqrt2}{2}
\]
因此
\[
S(a)\text{ 在 }\left(\frac12,\frac{1+\sqrt2}{2}\right)\text{ 上单调递增，}
\]
\[
S(a)\text{ 在 }\left(\frac{1+\sqrt2}{2},+\infty\right)\text{ 上单调递减.}
\]
并在 \(a=a_0\) 时取得最大值.

若取较大的最大值点 \(x_3=1-\frac1{4a}\)，则
\[
\widetilde S(a)
=\frac12(x_2-x_1)\left(1-\frac1{4a}\right)
=\frac{(2a-1)(4a-1)}{4(1+4a^2)}
\]
求导得
\(\widetilde S'(a) =\frac{2(12a^2+4a-3)}{4(1+4a^2)^2}>0\)，\(\left(a>\frac12\right)\)
所以此时 \(\widetilde S(a)\) 在 \((1/2,+\infty)\) 上单调递增. 由此可见，若不补充 \(x_3\) 取哪一个最大值点，第 3 问的 \(S(a)\) 并不能唯一确定；上面的两种结论分别对应两种合法选择.
\end{enumerate}
\end{solution}
